\documentclass[12pt,a4paper]{exam}
\usepackage[utf8]{inputenc}
\usepackage{amsmath,amssymb,amsfonts}
\usepackage{geometry}
\usepackage{enumitem}
\usepackage{graphicx}
\usepackage{hyperref}
\usepackage{xcolor}
\geometry{a4paper, top=2cm, bottom=2cm, left=2.5cm, right=2.5cm}
\hypersetup{colorlinks=true, linkcolor=blue, urlcolor=blue}
\renewcommand{\thequestion}{\textbf{\arabic{question}}}
\renewcommand{\questionlabel}{\thequestion.}
\pointname{ marks}
\pointsdroppedatright
\bracketedpoints
\setlength{\parindent}{0pt}
\setlength{\emergencystretch}{3em}
\allowdisplaybreaks
\newsavebox{\schmathbox}
\newcommand{\fitmath}[1]{%
  \sbox{\schmathbox}{$\displaystyle #1$}%
  \ifdim\wd\schmathbox>\linewidth
    \resizebox{\linewidth}{!}{\usebox{\schmathbox}}%
  \else
    \usebox{\schmathbox}%
  \fi}

\begin{document}

\thispagestyle{empty}
\begin{center}
  \textbf{\LARGE Diag} \\[0.3cm]
  \textbf{Time Allowed: 3 Hours} \hfill \textbf{Maximum Marks: 80} \\[0.3cm]
  \rule{\textwidth}{0.5pt}
\end{center}

\vspace{0.3cm}

\noindent\textbf{General Instructions:}
\begin{enumerate}
  \item {\normalfont\normalcolor This question paper contains 40 questions. All questions are compulsory.}
  \item {\normalfont\normalcolor Questions are grouped into sections A through E.}
  \item {\normalfont\normalcolor Use of calculators is NOT permitted.}
\end{enumerate}
\bigskip
\noindent\textbf{\large Section A: Multiple Choice \& Assertion-Reason (1 mark each)}

\begin{questions}
\question[1] Two poles of height $6$ m and $11$ m stand on a plane ground. If the distance between their feet is $12$ m, find the distance between their tops.
\begin{enumerate}[label=(\Alph*)]
  \item (A) 12 m
  \item (B) 13 m
  \item (C) 14 m
  \item (D) 15 m
\end{enumerate}
\vspace{0.15cm}

\question[1] In $\triangle ABC$, $D$ and $E$ are points on sides $AB$ and $AC$ respectively such that $AD = 3$ cm, $DB = 2$ cm, $AE = 4.5$ cm, and $EC = 3$ cm. Which of the following is true?
\begin{enumerate}[label=(\Alph*)]
  \item (A) $DE \parallel BC$
  \item (B) $DE \parallel AB$
  \item (C) $DE \parallel AC$
  \item (D) $DE$ is not parallel to any side
\end{enumerate}
\vspace{0.15cm}

\question[1] If $\triangle ABC \sim \triangle DEF$, $\angle A = 50^\circ$, and $\angle E = 70^\circ$, then $\angle C$ is:
\begin{enumerate}[label=(\Alph*)]
  \item $(A) 50^\circ$
  \item $(B) 70^\circ$
  \item $(C) 60^\circ$
  \item $(D) 80^\circ$
\end{enumerate}
\vspace{0.15cm}

\question[1] The discriminant of the quadratic equation $3x^2 - 4x + 1 = 0$ is
\begin{enumerate}[label=(\Alph*)]
  \item (A) 1
  \item (B) 2
  \item (C) 4
  \item (D) 16
\end{enumerate}
\vspace{0.15cm}

\question[1] The discriminant of the quadratic equation $3x^2 - 5x + 2 = 0$ is
\begin{enumerate}[label=(\Alph*)]
  \item (A) 25
  \item (B) 1
  \item (C) 49
  \item (D) -1
\end{enumerate}
\vspace{0.15cm}

\question[1] The discriminant of the quadratic equation $2x^2 - 5x + 3 = 0$ is
\begin{enumerate}[label=(\Alph*)]
  \item (A) 1
  \item (B) -1
  \item (C) 25
  \item (D) 49
\end{enumerate}
\vspace{0.15cm}

\question[1] If the prime factorization of a positive integer $n$ is given by $n = 2^3 \text{imes} 3^2 \text{imes} 5^k$ and it is known that $n$ is a multiple of $150$, then the smallest possible value of the non-negative integer $k$ is:
\begin{enumerate}[label=(\Alph*)]
  \item (A) 0
  \item (B) 1
  \item (C) 2
  \item (D) 3
\end{enumerate}
\vspace{0.15cm}

\question[1] The roots of the quadratic equation $x^2 - 7x + 12 = 0$ are obtained by factorising the middle term. The roots of the equation are:
\begin{enumerate}[label=(\Alph*)]
  \item (A) $2, 6$
  \item (B) $3, 4$
  \item (C) $-3, -4$
  \item (D) $1, 12$
\end{enumerate}
\vspace{0.15cm}

\question[1] For what value of $k$ does the quadratic equation $(k-2)x^2 + 2(2k-3)x + (5k-6) = 0$ have equal real roots?
\begin{enumerate}[label=(\Alph*)]
  \item (A) $3, 1$
  \item (B) $3, -1$
  \item (C) $-3, 1$
  \item (D) $2, 3$
\end{enumerate}
\vspace{0.15cm}

\question[1] If one root of the quadratic equation $x^2 - 8x + m = 0$ exceeds the other by $2$, then the value of $m$ is:
\begin{enumerate}[label=(\Alph*)]
  \item (A) $12$
  \item (B) $14$
  \item (C) $15$
  \item (D) $16$
\end{enumerate}
\vspace{0.15cm}

\question[1] If the prime factorization of a positive integer $n$ is given by $n = 2^3 \text{imes} 3^2 \text{imes p}$ and the prime factorization of $m$ is $2^2 \text{imes} 3 \text{imes q}$, where $p$ and $q$ are distinct prime numbers other than $2$ and $3$, then the HCF of $n$ and $m$ is:
\begin{enumerate}[label=(\Alph*)]
  \item (A) $12$
  \item (B) $6$
  \item (C) $24$
  \item (D) $72$
\end{enumerate}
\vspace{0.15cm}

\question[1] If the polynomial $p(x) = x^4 - 6x^3 + 16x^2 - 25x + 10$ is divided by $g(x) = x^2 - 2x + k$, the remainder is found to be $x + a$. The value of $k$ is:
\begin{enumerate}[label=(\Alph*)]
  \item (A) 2
  \item (B) 3
  \item (C) 5
  \item (D) 7
\end{enumerate}
\vspace{0.15cm}

\question[1] Assertion (A): If 	riangle ABC $\sim$ 	riangle DEF and $AB = 2$ cm, $DE = 4$ cm, then the ratio of their areas is $1:4$. Reason (R): The ratio of the areas of two similar triangles is equal to the square of the ratio of their corresponding sides.
\begin{enumerate}[label=(\Alph*)]
  \item (A) Both Assertion (A) and Reason (R) are true and Reason (R) is the correct explanation of Assertion (A)
  \item (B) Both Assertion (A) and Reason (R) are true but Reason (R) is NOT the correct explanation of Assertion (A)
  \item (C) Assertion (A) is true but Reason (R) is false
  \item (D) Assertion (A) is false but Reason (R) is true
\end{enumerate}
\vspace{0.15cm}

\question[1] Assertion (A): In $\text{riangle ABC}$, if $DE \parallel BC$ with $D$ on $AB$ and $E$ on $AC$, then $AD/DB = AE/EC$. Reason (R): Thales Theorem (Basic Proportionality Theorem) states that if a line is drawn parallel to one side of a triangle intersecting the other two sides, then it divides the two sides in the same ratio.

\begin{center}
\includegraphics[width=0.48\textwidth]{/home/user/Desktop/backup/schEngine/pdfs/diag/diagrams/diagram_105.pdf}
\end{center}

\begin{enumerate}[label=(\Alph*)]
  \item (A) Both Assertion (A) and Reason (R) are true and Reason (R) is the correct explanation of Assertion (A)
  \item (B) Both Assertion (A) and Reason (R) are true but Reason (R) is NOT the correct explanation of Assertion (A)
  \item (C) Assertion (A) is true but Reason (R) is false
  \item (D) Assertion (A) is false but Reason (R) is true
\end{enumerate}
\vspace{0.15cm}

\question[1] Assertion (A): Two triangles are similar if their corresponding angles are equal. Reason (R): If the corresponding sides of two triangles are in the same ratio, then the triangles are congruent.
\begin{enumerate}[label=(\Alph*)]
  \item (A) Both Assertion (A) and Reason (R) are true and Reason (R) is the correct explanation of Assertion (A)
  \item (B) Both Assertion (A) and Reason (R) are true but Reason (R) is NOT the correct explanation of Assertion (A)
  \item (C) Assertion (A) is true but Reason (R) is false
  \item (D) Assertion (A) is false but Reason (R) is true
\end{enumerate}
\vspace{0.15cm}

\question[1] Assertion (A): In $\text{riangle ABC}$, if $\angle A = 90^\circ$ and $AD \perp BC$, then 	riangle ADB $\sim$ 	riangle CDA. Reason (R): Two triangles are similar if two angles of one triangle are respectively equal to two angles of the other triangle.
\begin{enumerate}[label=(\Alph*)]
  \item (A) Both Assertion (A) and Reason (R) are true and Reason (R) is the correct explanation of Assertion (A)
  \item (B) Both Assertion (A) and Reason (R) are true but Reason (R) is NOT the correct explanation of Assertion (A)
  \item (C) Assertion (A) is true but Reason (R) is false
  \item (D) Assertion (A) is false but Reason (R) is true
\end{enumerate}
\vspace{0.15cm}

\question[1] Assertion (A): The perimeters of two similar triangles are in the same ratio as their corresponding sides. Reason (R): If 	riangle ABC $\sim$ 	riangle PQR, then $AB/PQ = BC/QR = CA/RP = (AB+BC+CA)/(PQ+QR+RP)$.
\begin{enumerate}[label=(\Alph*)]
  \item (A) Both Assertion (A) and Reason (R) are true and Reason (R) is the correct explanation of Assertion (A)
  \item (B) Both Assertion (A) and Reason (R) are true but Reason (R) is NOT the correct explanation of Assertion (A)
  \item (C) Assertion (A) is true but Reason (R) is false
  \item (D) Assertion (A) is false but Reason (R) is true
\end{enumerate}
\vspace{0.15cm}

\question[1] Assertion (A): For any two positive integers $a$ and $b$, if $HCF(a, b) = 11$ and $LCM(a, b) = 693$, then the product $a \times b = 7623$. Reason (R): For any two positive integers $a$ and $b$, $HCF(a, b) \times LCM(a, b) = a \times b$.
\begin{enumerate}[label=(\Alph*)]
  \item (A) Both Assertion (A) and Reason (R) are true and Reason (R) is the correct explanation of Assertion (A)
  \item (B) Both Assertion (A) and Reason (R) are true but Reason (R) is NOT the correct explanation of Assertion (A)
  \item (C) Assertion (A) is true but Reason (R) is false
  \item (D) Assertion (A) is false but Reason (R) is true
\end{enumerate}
\vspace{0.15cm}

\question[1] Assertion (A): For the quadratic equation $3x^2 - kx + 3 = 0$, if the roots are equal, then $k = 6$ or $k = -6$. \textbackslash{}\textbackslash{} Reason (R): A quadratic equation $ax^2 + bx + c = 0$ has equal real roots if and only if its discriminant $D = b^2 - 4ac$ is equal to $0$.
\begin{enumerate}[label=(\Alph*)]
  \item (A) Both Assertion (A) and Reason (R) are true and Reason (R) is the correct explanation of Assertion (A)
  \item (B) Both Assertion (A) and Reason (R) are true but Reason (R) is NOT the correct explanation of Assertion (A)
  \item (C) Assertion (A) is true but Reason (R) is false
  \item (D) Assertion (A) is false but Reason (R) is true
\end{enumerate}
\vspace{0.15cm}

\question[1] Assertion (A): The number $\frac{7}{2^3 \times 5^2}$ has a terminating decimal expansion.
Reason (R): A rational number $\frac{p}{q}$ has a terminating decimal expansion if the prime factorization of $q$ is of the form $2^n \times 5^m$, where $n$ and $m$ are non-negative integers.
\begin{enumerate}[label=(\Alph*)]
  \item (A) Both Assertion (A) and Reason (R) are true and Reason (R) is the correct explanation of Assertion (A)
  \item (B) Both Assertion (A) and Reason (R) are true but Reason (R) is NOT the correct explanation of Assertion (A)
  \item (C) Assertion (A) is true but Reason (R) is false
  \item (D) Assertion (A) is false but Reason (R) is true
\end{enumerate}
\vspace{0.15cm}

\end{questions}
\bigskip
\noindent\textbf{\large Section B: Very Short Answer (2 marks each)}

\begin{questions}
\question[2] In $\triangle ABC$, $DE \parallel BC$ such that $AD = 3$ $cm$, $DB = 4$ $cm$, and $AE = 6$ $cm$. Find the length of $EC$.

\begin{center}
\includegraphics[width=0.48\textwidth]{/home/user/Desktop/backup/schEngine/pdfs/diag/diagrams/diagram_106.pdf}
\end{center}

\vspace{0.15cm}

\question[2] Two triangles $\triangle ABC$ and $\triangle DEF$ are similar. If $\text{Area}(\triangle ABC) = 64$ $cm^2$ and $\text{Area}(\triangle DEF) = 121$ $cm^2$, and $EF = 11$ $cm$, find the length of $BC$.
\vspace{0.15cm}

\question[2] In $\triangle PQR$, $S$ and $T$ are points on sides $PQ$ and $PR$ respectively such that $\angle PST = \angle PQR$. If $PS = 2$ $cm$, $SQ = 3$ $cm$, and $PT = 1.5$ $cm$, find the length of $TR$.
\vspace{0.15cm}

\question[2] A vertical pole of length $6$ $m$ casts a shadow $4$ $m$ long on the ground and at the same time a tower casts a shadow $28$ $m$ long. Find the height of the tower.
\vspace{0.15cm}

\question[2] In $\triangle ABC$, $AD \perp BC$. Prove that $AB^2 + CD^2 = AC^2 + BD^2$.

\begin{center}
\includegraphics[width=0.48\textwidth]{/home/user/Desktop/backup/schEngine/pdfs/diag/diagrams/diagram_107.pdf}
\end{center}

\vspace{0.15cm}

\question[2] Find the value of $k$ for which the quadratic equation $(k-1)x^2 - 2(k-1)x + 1 = 0$ has equal real roots.
\vspace{0.15cm}

\question[2] If the HCF of $126$ and $45$ is expressed in the form $126m + 45n$, where $m$ and $n$ are integers, find one possible pair of values for $(m, n)$ such that $m + n$ is minimized.
\vspace{0.15cm}

\end{questions}
\bigskip
\noindent\textbf{\large Section C: Short Answer (3 marks each)}

\begin{questions}
\question[3] In $\triangle ABC$, $DE \parallel BC$ such that $AD = (4x - 3) \text{ cm}$, $AE = (8x - 7) \text{ cm}$, $BD = (3x - 1) \text{ cm}$, and $CE = (5x - 3) \text{ cm}$. Find the value of $x$.

\begin{center}
\includegraphics[width=0.48\textwidth]{/home/user/Desktop/backup/schEngine/pdfs/diag/diagrams/diagram_108.pdf}
\end{center}

\vspace{0.15cm}

\question[3] Two poles of height $6 \text{ m}$ and $11 \text{ m}$ stand on a plane ground. If the distance between their feet is $12 \text{ m}$, find the distance between their tops.
\vspace{0.15cm}

\question[3] In $\triangle ABC$, $\angle A = 90^\circ$ and $AD \perp BC$. If $AD = 8 \text{ cm}$ and $BD = 4 \text{ cm}$, find the length of $CD$.
\vspace{0.15cm}

\question[3] The perimeters of two similar triangles are $30 \text{ cm}$ and $20 \text{ cm}$ respectively. If one side of the first triangle is $9 \text{ cm}$, find the length of the corresponding side of the second triangle.
\vspace{0.15cm}

\question[3] In $\triangle ABC$, $P$ and $Q$ are points on sides $AB$ and $AC$ such that $PQ \parallel BC$. If $AP = 3 \text{ cm}$, $PB = 2 \text{ cm}$ and $AQ = 2.7 \text{ cm}$, find $AC$.

\begin{center}

\includegraphics[width=0.48\textwidth]{/home/user/Desktop/backup/schEngine/pdfs/diag/diagrams/diagram_110.pdf}
\end{center}

\vspace{0.15cm}

\question[3] If $\alpha$ and $\beta$ are the roots of the quadratic equation $x^2 - 5x + 3 = 0$, find the value of $\frac{1}{\alpha} + \frac{1}{\beta} - \alpha\beta$.
\vspace{0.15cm}

\end{questions}
\bigskip
\noindent\textbf{\large Section D: Long Answer (4-5 marks each)}

\begin{questions}
\question[5] Prove that if a line is drawn parallel to one side of a triangle to intersect the other two sides in distinct points, the other two sides are divided in the same ratio. Use this theorem to find the length of $EC$ in $\text{riangle ABC}$ where $DE \parallel BC$, $AD = 3$ cm, $DB = 4$ cm, and $AE = 6$ cm.

\begin{center}
\includegraphics[width=0.48\textwidth]{/home/user/Desktop/backup/schEngine/pdfs/diag/diagrams/diagram_111.pdf}
\end{center}

\vspace{0.15cm}

\question[5] In $\triangle ABC$, $\angle A = 90^\circ$ and $AD \perp BC$. Prove that $AD^2 = BD \times CD$. If $BD = 4$ cm and $CD = 9$ cm, calculate the length of $AD$.
\vspace{0.15cm}

\question[5] Two poles of heights $a$ meters and $b$ meters are $p$ meters apart. Prove that the height of the point of intersection of the lines joining the top of each pole to the foot of the opposite pole is given by $\frac{ab}{a+b}$ meters.
\vspace{0.15cm}

\question[5] Prove that the ratio of the areas of two similar triangles is equal to the square of the ratio of their corresponding sides. Using this, find the ratio of the areas of two similar triangles if their sides are in the ratio $4:9$.
\vspace{0.15cm}

\question[4] Show that for any positive integer $n$, $n^3 - n$ is divisible by $6$.
\vspace{0.15cm}

\question[4] Find the greatest number of 6 digits which is exactly divisible by $12, 15,$ and $20$.
\vspace{0.15cm}

\question[4] A merchant has $120$ litres of oil of one kind, $180$ litres of another kind and $240$ litres of a third kind. He wants to sell the oil by filling the three kinds of oil in tins of equal capacity. What should be the greatest capacity of such a tin?
\vspace{0.15cm}

\end{questions}
\bigskip

\end{document}
